\documentclass[11pt,a4paper]{article}
\usepackage[utf8]{inputenc}
\usepackage[T1]{fontenc}
\usepackage[french]{babel}
\usepackage{geometry}
\usepackage{xcolor}
\usepackage{hyperref}

\geometry{margin=1in}
\definecolor{lorraine_blue}{RGB}{0, 75, 142}

\begin{document}

\noindent \textbf{Ismail AISSAOUI} \\
Ingénieur d'État en Intelligence Artificielle (ESI-SBA) \\
Email : ismail.aissaoui.pro@gmail.com \\
Téléphone : +213 660 70 77 96 \\
LinkedIn : \href{https://linkedin.com/in/aissaoui-ismail-6a77a92a7}{linkedin.com/in/aissaoui-ismail-6a77a92a7}

\bigskip

\begin{flushright}
    À l'attention du Pr. Enrique Espinosa et du Pr. Sabeur Aridhi \\
    \textbf{Université de Lorraine} \\
    Laboratoires CRM2 / LORIA \\
    Nancy, France \\
    \medskip
    Le 11 mai 2026
\end{flushright}

\bigskip

\noindent \textbf{Objet : Candidature au doctorat : Prédiction de l'assemblage moléculaire dans les APIs par IA-Machine Learning}

\bigskip

Monsieur le Professeur Espinosa, Monsieur le Professeur Aridhi, Monsieur Aubert,

\medskip

C'est avec un vif intérêt que je soumets ma candidature pour le poste de doctorat portant sur la prédiction de l'assemblage moléculaire dans les principes actifs pharmaceutiques (APIs) via l'IA. Diplômé en tant qu'Ingénieur d'État en IA de l'ESI-SBA, avec une spécialisation en **Réseaux de neurones informés par la physique (PINN)** et en **Apprentissage sur Graphes**, je suis particulièrement motivé par le défi de lier les données quantiques à l'organisation 3D moléculaire.

Le caractère interdisciplinaire de ce projet, alliant l'expertise cristallographique du **CRM2** à l'excellence en IA du **LORIA**, correspond parfaitement à mes aspirations de recherche. Mon profil s'aligne sur les exigences de la thèse dans plusieurs domaines clés :

\begin{itemize}
    \item \textbf{Synergie Physique-IA} : Dans ma thèse d'ingénieur chez **Sonatrach**, je développe un **Jumeau Numérique Génératif** pour des turbines à gaz. Mes recherches se concentrent sur l'utilisation des **PINN** pour intégrer des lois physiques directement dans le processus d'apprentissage. Cette méthodologie est directement transposable à votre objectif d'apprendre des représentations à partir de **données quantiques** ; je comprends comment garantir que les prédictions de l'IA respectent les règles physico-chimiques sous-jacentes.
    \item \textbf{Modélisation Structurelle et Graphes} : La prédiction de l'assemblage moléculaire est un défi structurel multi-dimensionnel. Mon expertise en **Réseaux de Neurones sur Graphes (GNN)** et en architectures séquentielles avancées (**Mamba-SSM**, **xLSTM**) me dote des outils nécessaires pour modéliser les dépendances inter-atomiques complexes.
    \item \textbf{Gestion de Données Scientifiques} : J'ai l'expérience de la mise en place de pipelines ML pour des datasets scientifiques massifs (5,2 millions de points). Je maîtrise **PyTorch** et **Scikit-learn**, et je suis impatient d'appliquer ces compétences à votre base de données de 455 000 entrées issues de calculs DFT.
\end{itemize}

Au-delà de l'adéquation technique, je suis passionné par la philosophie du "Scientifique-dans-la-boucle", où l'IA sert à accélérer les découvertes en sciences physiques. Je suis une personne qui apprend vite et je me réjouis d'approfondir mes connaissances en chimie quantique sous votre direction.

Je suis convaincu que ma vision d'ingénieur en IA, combinée à ma rigueur académique, me permettra de contribuer de manière significative à l'objectif du projet d'anticiper les assemblages 3D des principes pharmaceutiques.

Je vous remercie de l'attention portée à ma candidature et reste à votre entière disposition pour un entretien.

\bigskip

\noindent Je vous prie d'agréer, Messieurs, l'expression de mes salutations distinguées.

\bigskip

\noindent Ismail AISSAOUI

\end{document}
